\documentclass[12pt]{article}
\usepackage{amsmath}

\usepackage{xcolor}

\usepackage{listings}
\usepackage[most]{tcolorbox}

% style stolen from haskell notes (non avrei ricreato appositamente altrimenti...)
\definecolor{mzbg}{HTML}{F8F9FA}
\definecolor{mzpurple}{HTML}{5E5086}
\definecolor{mzgreen}{HTML}{118C4F}
\definecolor{mzred}{HTML}{C0392B}
\definecolor{mzblue}{HTML}{2E86C1} 
\definecolor{mznumber}{HTML}{ADB5BD}
\definecolor{mztitlebg}{HTML}{2C3E50}

\lstdefinelanguage{MiniZinc}{
    morekeywords={include, var, par, array, of, int, float, set, bool, constraint, solve, satisfy, minimize, maximize, forall, exists, where, if, then, else, endif},
    morekeywords=[2]{alldifferent, abs, sum, count, min, max, card}, 
    sensitive=true,
    morecomment=[l]{\%},
    morecomment=[s]{/*}{*/},
    morestring=[b]",
}

\lstdefinestyle{mznstyle}{
    language=MiniZinc,
    commentstyle=\color{mzgreen}\itshape,
    keywordstyle=\color{mzpurple}\bfseries,
    keywordstyle=[2]\color{mzblue},
    numberstyle=\tiny\color{mznumber},
    stringstyle=\color{mzred},
    basicstyle=\ttfamily\footnotesize,
    breaklines=true,
    showstringspaces=false,
    tabsize=2,
    numbers=none
}

\newtcblisting{mznbox}[1]{
    enhanced,
    skin=bicolor,
    colback=mzbg,
    colbacktitle=mztitlebg,
    coltitle=white,
    colframe=mztitlebg,
    boxrule=0.5pt,
    arc=4pt,
    drop shadow=black!10!white,
    fonttitle=\bfseries\sffamily,
    title=#1,
    listing only,
    listing options={style=mznstyle, numbers=left, numbersep=10pt},
    left=15pt,
    top=4pt,
    bottom=4pt,
    breakable
}

\title{E.A.5.2 - CSP: edge colouring}
\author{}
\date{}

\begin{document}
\maketitle
\vspace{-5em}
\section{Modellazione}

Per evitare di creare due variabili per ogni arco (indiretto), impongo un ordinamento.

\begin{itemize}
    \item variabili: $X = \{ X_{a, b} \mid (a,b)\in E, \ a < b \}$
    \item domini: $D_{X_{a,b}} = \{1, 2, 3\} \ \forall X_{a,b}\in X$
    \item vincoli: 
        \begin{align*}
            C = &\{\neg\bigl( (X_{a,b} = X_{b,c}) \land (X_{b,c} = X_{a,c})\bigr) \\
                &\mid a, b, c \in V \text{ s.t. } a < b < c, \ \land \  (a,b), (b,c), (a,c) \in E\} 
        \end{align*}
\end{itemize}

\section{Istanziazione}

\subsection{Variabili e domini}
(Ordino i nodi alfabeticamente per rispettare la regola $a < b$ ed evitare duplicati)
\begin{align*}
X = \{ & X_{A,E}, X_{A,H}, X_{A,I}, X_{A,S}, \\
& X_{B,C}, X_{B,G2}, X_{B,I}, X_{B,J}, X_{B,S}, \\
& X_{C,D}, X_{C,G2}, X_{C,S}, \\
& X_{D,E}, X_{D,S}, \\
& X_{E,G1}, X_{E,H}, \\
& X_{G1,H}, \\
& X_{G2,J}, \\
& X_{H,I} \}
\end{align*}

Tutte le variabili condividono lo stesso dominio di colori (1, 2 o 3).
\[
D_v = \{1, 2, 3\} \quad \forall v \in X
\]

\subsection{Vincoli}
Ci sono 7 triangoli nel grafo:

\begin{itemize}
\item \textbf{\{A, E, H\}:}
        \[ C_1: \neg ((X_{A,E} = X_{A,H}) \land (X_{A,H} = X_{E,H})) \]

    \item \textbf{\{A, H, I\}:}
        \[ C_2: \neg ((X_{A,H} = X_{A,I}) \land (X_{A,I} = X_{H,I})) \]

    \item \textbf{\{B, C, G2\}:}
        \[ C_3: \neg ((X_{B,C} = X_{B,G2}) \land (X_{B,G2} = X_{C,G2})) \]

    \item \textbf{\{B, C, S\}:}
        \[ C_4: \neg ((X_{B,C} = X_{B,S}) \land (X_{B,S} = X_{C,S})) \]

    \item \textbf{\{B, G2, J\}:}
        \[ C_5: \neg ((X_{B,G2} = X_{B,J}) \land (X_{B,J} = X_{G2,J})) \]

    \item \textbf{\{C, D, S\}:}
        \[ C_6: \neg ((X_{C,D} = X_{C,S}) \land (X_{C,S} = X_{D,S})) \]

    \item \textbf{\{E, G1, H\}:}
        \[ C_7: \neg ((X_{E,G1} = X_{E,H}) \land (X_{E,H} = X_{G1,H})) \]
    \item \textbf{insieme finale:}
        \[
            C = \{ C_1, C_2, C_3, C_4, C_5, C_6, C_7 \}
        \]

\end{itemize}

\section{Codifica in MiniZinc}

\subsection{v.1: matrice di adiacenza}
Versione senza assunzioni con matrice di adiacenza.

\begin{mznbox}{}
int: N; 

% boolean matrix where true = edge
array[1..N, 1..N] of bool: E;

% variables: 0 = no edge, 1..3 are the colors
array[1..N, 1..N] of var 0..3: color;

% graph structure constraints
constraint forall(i, j in 1..N where i < j) (
    % undirected graph symmetry
    color[i, j] == color[j, i] /\ 

    % color iff edge
    (E[i, j] <-> color[i, j] > 0)
);

% no self-loops
constraint forall(i in 1..N) (color[i, i] == 0);

% triangle constraint
constraint forall(i, j, k in 1..N where i < j /\ j < k) (
    E[i, j] /\ E[j, k] /\ E[i, k] ->
        not (color[i, j] == color[j, k] /\ color[j, k] == color[i, k])
);

solve satisfy;  
\end{mznbox}


\subsection{v.2: in funzione degli archi}
Versione ``edge-centric''. Se facciamo anche assunzioni sui dati in input (che per ogni arco $(u, v)$ valga $u < v$, e che le liste dei nodi siano fornite in ordine crescente), possiamo identificare i triangoli in modo semplice.

Questa versione è più `pulita' a livello di modellazione (più vicina alla modellazione matematica), ma potenzialmente più lenta computazionalmente (nella pratica, non ho in realtà rilevato molta differenza a livello di tempo di risoluzione).

\begin{itemize}
    \item inserendo il \texttt{u[i] == u[j] /\ v[i] == u[k] /\ v[j] == v[k]} nel \texttt{where}, i triangoli vengono identificati a tempo di compilazione (e non come vincoli dal solver), rendendo il codice più performante (su grafo completo $K_{10}$, $\sim$50ms in meno)
\end{itemize}

\begin{mznbox}{}
int: N; % nodes
int: M; % edges

% edges: edge i = (u[i], v[i])
array[1..M] of int: u;
array[1..M] of int: v;

% we have to assign colors to edges 
array[1..M] of var 1..3: color;

% we can check that our assumptions hold:
constraint assert(forall(i in 1..M) (u[i] < v[i]), 
    "!(u[i] < v[i])");

constraint assert(forall(i in 1..M-1) (
    u[i] < u[i+1] \/ (u[i] == u[i+1] /\ v[i] < v[i+1])
), "u, v aren't ordered");

% if our assumptions hold, we can check:
constraint forall(i, j, k in 1..M where i < j /\ j < k /\ u[i] == u[j] /\ v[i] == u[k] /\ v[j] == v[k]) 
(
    not (color[i] == color[j] /\ color[j] == color[k])
);

solve satisfy;

\end{mznbox}

\end{document}
